%% kapa: exam=1 round=28 year=2017 period=2 form=A subject=appraisal_law subject_ko=감정평가관계법규 seq=29 number=29
%% source: Q-Net 감정평가사 기출문제 (cst00309) artlSeq=5199603 file=extracted/1008/5199603/2017년 제28회 감정평가사 1차시험 문제지/2교시 A형.pdf
%% conversion: kordoc 4.24.51 -> tools/gen1.py (auto); review=proofread-sonnet; flags=none
\begin{question}{29}
공간정보의 구축 및 관리 등에 관한 법령상 용어에 관한 설명으로 옳지 않은 것은?
\begin{choices}
  \item 지적측량은 지적확정측량 및 지적재조사측량을 포함한다.
  \item 필지를 구획하는 선의 굴곡점으로서 지적도에 도해(圖解) 형태로 등록하는 점은 경계점에 해당한다.
  \item 지적공부는 지적측량 등을 통하여 조사된 토지의 표시와 해당 토지의 소유자 등을 기록한 대장 및 도면을 말한다.
  \item 축척변경은 지적도에 등록된 경계점의 정밀도를 높이기 위하여 작은 축척을 큰 축척으로 변경하여 등록하는 것을 말한다.
  \item 등록전환은 토지대장 및 지적도에 등록된 임야를 임야대장 및 임야도에 옮겨 등록하는 것을 말한다.
\end{choices}
\end{question}
\answer{5}
